\documentclass[a4paper,11pt]{article}

\usepackage[a4paper,margin=2cm]{geometry}
\usepackage[T1]{fontenc}
\usepackage[utf8]{inputenc}
\usepackage[ngerman,english]{babel}
\usepackage{lmodern}
\usepackage{microtype}
\usepackage[table]{xcolor}
\usepackage{parskip}
\usepackage{enumitem}
\usepackage[most]{tcolorbox}
\usepackage{titlesec}
\usepackage[hidelinks]{hyperref}
\usepackage{array}

\definecolor{HeaderBlueDark}{RGB}{70,110,170}
\definecolor{RowBlueLight}{RGB}{235,243,255}
\definecolor{RowBlueDark}{RGB}{220,233,250}
\definecolor{SoftGray}{RGB}{90,90,90}

\setlength{\parindent}{0pt}
\setlist[itemize]{leftmargin=1.4em,itemsep=0.35em,topsep=0.35em}
\linespread{1.06}

\titleformat{\section}[block]
{\Large\bfseries\color{HeaderBlueDark}}
{}{0pt}{}
\titlespacing{\section}{0pt}{1.3em}{0.8em}

\titleformat{\subsection}[block]
{\large\bfseries\color{HeaderBlueDark}}
{}{0pt}{}
\titlespacing{\subsection}{0pt}{1.0em}{0.6em}

\newtcolorbox{exbox}{enhanced,breakable,colback=RowBlueLight,colframe=RowBlueDark,boxrule=0.4pt,arc=1.5mm,left=1.2mm,right=1.2mm,top=1mm,bottom=1mm,title={Beispiele \textnormal{(Examples)}},fonttitle=\bfseries,colbacktitle=RowBlueDark,coltitle=black}

\NewDocumentEnvironment{grammarentry}{mm+b}{%
    \begin{tcolorbox}[enhanced,breakable,colback=white,colframe=HeaderBlueDark,boxrule=0.6pt,arc=1.5mm,left=1.5mm,right=1.5mm,top=1mm,bottom=1mm,colbacktitle=HeaderBlueDark,coltitle=white,fonttitle=\bfseries,title={#1\hfill\normalfont\footnotesize #2}]
    #3
    \end{tcolorbox}
}{}

\newcommand{\GermanText}[1]{\textbf{Deutsch: }#1\par}
\newcommand{\EnglishText}[1]{{\small\color{SoftGray}\textbf{English: }#1\par}}
\newcommand{\ExampleItem}[2]{\item \textbf{#1}\\{\small\color{SoftGray}#2}}

\title{Starke, schwache und gemischte Adjektivendungen}
\date{}

\begin{document}
    \maketitle
    \tableofcontents
    \newpage

\section{Grundidee}

\begin{grammarentry}{Warum gibt es starke und schwache Endungen?}{Why strong and weak endings exist}
\GermanText{Das grammatische Thema heißt \textbf{Adjektivdeklination}. Es geht darum, welche Endung ein Adjektiv bekommt, wenn es vor einem Nomen steht.}
\EnglishText{The grammar topic is called \textbf{adjective declension}. It is about which ending an adjective gets when it stands before a noun.}

\GermanText{Die wichtigste Idee ist: Im Deutschen muss die Nomengruppe zeigen, welches \textbf{Genus}, welcher \textbf{Kasus} und welcher \textbf{Numerus} gemeint ist.}
\EnglishText{The most important idea is: In German, the noun phrase must show which \textbf{gender}, \textbf{case} and \textbf{number} are meant.}

\GermanText{Diese grammatische Information kann entweder vom Artikel oder vom Adjektiv getragen werden.}
\EnglishText{This grammatical information can be carried either by the article or by the adjective.}

\begin{exbox}
\begin{itemize}
    \ExampleItem{mit klarem Bezugswort}{with a clear antecedent}
    \ExampleItem{mit einem klaren Bezugswort}{with a clear antecedent}
    \ExampleItem{mit dem klaren Bezugswort}{with the clear antecedent}
\end{itemize}
\end{exbox}

\GermanText{In \textbf{mit klarem Bezugswort} gibt es keinen Artikel. Deshalb muss das Adjektiv \textbf{klar} die starke Endung tragen: \textbf{klarem}.}
\EnglishText{In \textbf{mit klarem Bezugswort}, there is no article. Therefore the adjective \textbf{klar} must carry the strong ending: \textbf{klarem}.}

\GermanText{In \textbf{mit einem klaren Bezugswort} trägt \textbf{einem} schon die wichtige Dativ-Neutrum-Information. Deshalb bekommt das Adjektiv nur die schwächere Endung: \textbf{klaren}.}
\EnglishText{In \textbf{mit einem klaren Bezugswort}, \textbf{einem} already carries the important dative-neuter information. Therefore the adjective only gets the weaker ending: \textbf{klaren}.}
\end{grammarentry}

\section{Die drei Haupttypen}

\begin{grammarentry}{Überblick}{overview}
\GermanText{Es gibt drei wichtige Muster: \textbf{starke Deklination}, \textbf{schwache Deklination} und \textbf{gemischte Deklination}.}
\EnglishText{There are three important patterns: \textbf{strong declension}, \textbf{weak declension} and \textbf{mixed declension}.}

\begin{center}
\rowcolors{2}{RowBlueLight}{RowBlueDark}
\begin{tabular}{>{\bfseries}p{3.5cm}p{5cm}p{5cm}}
\rowcolor{HeaderBlueDark}
\color{white}\textbf{Typ} & \color{white}\textbf{Wann?} & \color{white}\textbf{Beispiel} \\
stark & ohne Artikel & mit klarem Bezugswort \\
schwach & mit bestimmtem Artikel & mit dem klaren Bezugswort \\
gemischt & mit ein, kein, mein, dein usw. & mit einem klaren Bezugswort \\
\end{tabular}
\end{center}

\GermanText{Merksatz: Wenn der Artikel die Information klar zeigt, ist die Adjektivendung schwächer. Wenn kein Artikel da ist, muss das Adjektiv stärker arbeiten.}
\EnglishText{Memory rule: If the article clearly shows the grammatical information, the adjective ending is weaker. If there is no article, the adjective must work harder.}
\end{grammarentry}

\section{Starke Adjektivendungen}

\begin{grammarentry}{Starke Deklination}{strong adjective endings}
\GermanText{Die starke Deklination benutzt man meistens, wenn kein Artikel vor dem Adjektiv steht. Dann trägt das Adjektiv die wichtige grammatische Information.}
\EnglishText{Strong declension is usually used when there is no article before the adjective. Then the adjective carries the important grammatical information.}

\begin{exbox}
\begin{itemize}
    \ExampleItem{guter Kaffee}{good coffee}
    \ExampleItem{kaltes Wasser}{cold water}
    \ExampleItem{mit gutem Kaffee}{with good coffee}
    \ExampleItem{wegen schlechten Wetters}{because of bad weather}
\end{itemize}
\end{exbox}

\GermanText{Bei starken Endungen sieht man oft Endungen, die den bestimmten Artikeln ähnlich sind: \textbf{der} $\rightarrow$ \textbf{-er}, \textbf{das} $\rightarrow$ \textbf{-es}, \textbf{dem} $\rightarrow$ \textbf{-em}.}
\EnglishText{With strong endings, you often see endings that are similar to definite articles: \textbf{der} $\rightarrow$ \textbf{-er}, \textbf{das} $\rightarrow$ \textbf{-es}, \textbf{dem} $\rightarrow$ \textbf{-em}.}
\end{grammarentry}

\begin{center}
\rowcolors{2}{RowBlueLight}{RowBlueDark}
\begin{tabular}{>{\bfseries}p{2.5cm}p{3cm}p{3cm}p{3cm}p{3cm}}
\rowcolor{HeaderBlueDark}
\color{white}\textbf{Kasus} & \color{white}\textbf{maskulin} & \color{white}\textbf{feminin} & \color{white}\textbf{neutral} & \color{white}\textbf{Plural} \\
Nominativ & guter Mann & gute Frau & gutes Kind & gute Kinder \\
Akkusativ & guten Mann & gute Frau & gutes Kind & gute Kinder \\
Dativ & gutem Mann & guter Frau & gutem Kind & guten Kindern \\
Genitiv & guten Mannes & guter Frau & guten Kindes & guter Kinder \\
\end{tabular}
\end{center}

\section{Schwache Adjektivendungen}

\begin{grammarentry}{Schwache Deklination}{weak adjective endings}
\GermanText{Die schwache Deklination benutzt man nach bestimmten Artikeln und ähnlichen Wörtern. Der Artikel zeigt die Grammatik schon deutlich. Deshalb bekommt das Adjektiv meistens nur \textbf{-e} oder \textbf{-en}.}
\EnglishText{Weak declension is used after definite articles and similar words. The article already clearly shows the grammar. Therefore the adjective usually only gets \textbf{-e} or \textbf{-en}.}

\GermanText{Typische Wörter: \textbf{der, die, das, den, dem, des, dieser, jener, jeder, welcher, alle}.}
\EnglishText{Typical words: \textbf{der, die, das, den, dem, des, this, that, every, which, all}.}

\begin{exbox}
\begin{itemize}
    \ExampleItem{der gute Mann}{the good man}
    \ExampleItem{die gute Frau}{the good woman}
    \ExampleItem{das gute Kind}{the good child}
    \ExampleItem{mit dem guten Mann}{with the good man}
\end{itemize}
\end{exbox}

\GermanText{Die einfache Regel für schwache Endungen: \textbf{-e} im Nominativ Singular und im Akkusativ Feminin/Neutral. Fast überall sonst: \textbf{-en}.}
\EnglishText{The simple rule for weak endings: \textbf{-e} in nominative singular and in accusative feminine/neuter. Almost everywhere else: \textbf{-en}.}
\end{grammarentry}

\begin{center}
\rowcolors{2}{RowBlueLight}{RowBlueDark}
\begin{tabular}{>{\bfseries}p{2.5cm}p{3cm}p{3cm}p{3cm}p{3cm}}
\rowcolor{HeaderBlueDark}
\color{white}\textbf{Kasus} & \color{white}\textbf{maskulin} & \color{white}\textbf{feminin} & \color{white}\textbf{neutral} & \color{white}\textbf{Plural} \\
Nominativ & der gute Mann & die gute Frau & das gute Kind & die guten Kinder \\
Akkusativ & den guten Mann & die gute Frau & das gute Kind & die guten Kinder \\
Dativ & dem guten Mann & der guten Frau & dem guten Kind & den guten Kindern \\
Genitiv & des guten Mannes & der guten Frau & des guten Kindes & der guten Kinder \\
\end{tabular}
\end{center}

\section{Gemischte Adjektivendungen}

\begin{grammarentry}{Gemischte Deklination}{mixed adjective endings}
\GermanText{Die gemischte Deklination benutzt man nach \textbf{ein, kein, mein, dein, sein, ihr, unser, euer}.}
\EnglishText{Mixed declension is used after \textbf{ein, kein, mein, dein, sein, ihr, unser, euer}.}

\GermanText{Sie heißt gemischt, weil der Artikel manchmal die grammatische Information zeigt und manchmal nicht. Wenn der Artikel keine klare Endung hat, muss das Adjektiv stärker sein.}
\EnglishText{It is called mixed because the article sometimes shows the grammatical information and sometimes does not. If the article has no clear ending, the adjective must be stronger.}

\begin{exbox}
\begin{itemize}
    \ExampleItem{ein guter Mann}{a good man}
    \ExampleItem{ein gutes Kind}{a good child}
    \ExampleItem{eine gute Frau}{a good woman}
    \ExampleItem{mit einem guten Mann}{with a good man}
\end{itemize}
\end{exbox}

\GermanText{Bei \textbf{ein guter Mann} hat \textbf{ein} keine klare maskuline Nominativ-Endung. Deshalb bekommt das Adjektiv \textbf{-er}: \textbf{guter}.}
\EnglishText{In \textbf{ein guter Mann}, \textbf{ein} has no clear masculine nominative ending. Therefore the adjective gets \textbf{-er}: \textbf{guter}.}

\GermanText{Bei \textbf{mit einem guten Mann} zeigt \textbf{einem} schon Dativ Maskulin. Deshalb bekommt das Adjektiv nur \textbf{-en}: \textbf{guten}.}
\EnglishText{In \textbf{mit einem guten Mann}, \textbf{einem} already shows dative masculine. Therefore the adjective only gets \textbf{-en}: \textbf{guten}.}
\end{grammarentry}

\begin{center}
\rowcolors{2}{RowBlueLight}{RowBlueDark}
\begin{tabular}{>{\bfseries}p{2.5cm}p{3cm}p{3cm}p{3cm}p{3cm}}
\rowcolor{HeaderBlueDark}
\color{white}\textbf{Kasus} & \color{white}\textbf{maskulin} & \color{white}\textbf{feminin} & \color{white}\textbf{neutral} & \color{white}\textbf{Plural mit kein} \\
Nominativ & ein guter Mann & eine gute Frau & ein gutes Kind & keine guten Kinder \\
Akkusativ & einen guten Mann & eine gute Frau & ein gutes Kind & keine guten Kinder \\
Dativ & einem guten Mann & einer guten Frau & einem guten Kind & keinen guten Kindern \\
Genitiv & eines guten Mannes & einer guten Frau & eines guten Kindes & keiner guten Kinder \\
\end{tabular}
\end{center}

\section{Warum heißt es: mit einem klaren Bezugswort?}

\begin{grammarentry}{Analyse des Beispiels}{analysis of the example}
\GermanText{Der Ausdruck ist: \textbf{mit einem klaren Bezugswort}.}
\EnglishText{The expression is: \textbf{mit einem klaren Bezugswort}.}

\GermanText{\textbf{mit} verlangt Dativ. \textbf{das Bezugswort} ist neutral. Dativ Neutral mit \textbf{ein} ist \textbf{einem}.}
\EnglishText{\textbf{mit} requires dative. \textbf{das Bezugswort} is neuter. Dative neuter with \textbf{ein} is \textbf{einem}.}

\GermanText{Weil \textbf{einem} schon die starke Dativ-Endung \textbf{-em} trägt, braucht das Adjektiv nicht noch einmal \textbf{-em}. Deshalb heißt es \textbf{klaren}, nicht \textbf{klarem}.}
\EnglishText{Because \textbf{einem} already carries the strong dative ending \textbf{-em}, the adjective does not need \textbf{-em} again. Therefore it is \textbf{klaren}, not \textbf{klarem}.}

\begin{exbox}
\begin{itemize}
    \ExampleItem{mit klarem Bezugswort}{with a clear antecedent; no article, so the adjective is strong}
    \ExampleItem{mit einem klaren Bezugswort}{with a clear antecedent; mixed declension after \textbf{ein}}
    \ExampleItem{mit dem klaren Bezugswort}{with the clear antecedent; weak declension after \textbf{dem}}
\end{itemize}
\end{exbox}

\GermanText{Die falsche Form \textbf{mit einem klarem Bezugswort} ist eine Mischung aus zwei Signalen: \textbf{einem} hat schon \textbf{-em}, und \textbf{klarem} hätte noch einmal \textbf{-em}. Das ist in normalem Deutsch falsch.}
\EnglishText{The wrong form \textbf{mit einem klarem Bezugswort} is a mixture of two signals: \textbf{einem} already has \textbf{-em}, and \textbf{klarem} would have \textbf{-em} again. This is wrong in normal German.}
\end{grammarentry}

\section{Wichtige Wörter vor Adjektiven}

\begin{grammarentry}{Welche Wörter machen welchen Typ?}{which words cause which pattern}
\GermanText{Diese Übersicht hilft dir zu entscheiden, welches Muster du brauchst.}
\EnglishText{This overview helps you decide which pattern you need.}

\begin{center}
\rowcolors{2}{RowBlueLight}{RowBlueDark}
\begin{tabular}{>{\bfseries}p{4cm}p{4cm}p{6cm}}
\rowcolor{HeaderBlueDark}
\color{white}\textbf{Wort vor dem Adjektiv} & \color{white}\textbf{Typ} & \color{white}\textbf{Beispiel} \\
kein Artikel & stark & mit gutem Deutsch \\
der, die, das, dem, den, des & schwach & mit dem guten Deutsch \\
dieser, jener, jeder, welcher & schwach & mit diesem guten Beispiel \\
ein, eine, einem, einen, eines & gemischt & mit einem guten Beispiel \\
mein, dein, sein, ihr, unser, euer & gemischt & mit meinem guten Freund \\
kein & gemischt & mit keinem guten Grund \\
viele, einige, mehrere, wenige & oft stark oder wie Artikelwort, je nach Gebrauch & viele gute Beispiele \\
\end{tabular}
\end{center}
\end{grammarentry}

\section{Das wichtigste Muster für B1}

\begin{grammarentry}{Praktische Regel}{practical rule}
\GermanText{Für B1 kannst du dir diese einfache Reihenfolge merken:}
\EnglishText{For B1, you can remember this simple sequence:}

\begin{enumerate}
    \item \textbf{Gibt es keinen Artikel?} Dann meistens starke Endung: \textbf{mit gutem Deutsch}.
    \item \textbf{Gibt es der, die, das, dem, den, des?} Dann schwache Endung: \textbf{mit dem guten Deutsch}.
    \item \textbf{Gibt es ein, kein, mein, dein usw.?} Dann gemischte Endung: \textbf{mit einem guten Beispiel}.
\end{enumerate}

\EnglishText{1. If there is no article, usually use the strong ending. 2. If there is a definite article, use the weak ending. 3. If there is ein/kein/my/your etc., use the mixed ending.}

\begin{exbox}
\begin{itemize}
    \ExampleItem{Ich trinke kalten Kaffee.}{I drink cold coffee.}
    \ExampleItem{Ich trinke den kalten Kaffee.}{I drink the cold coffee.}
    \ExampleItem{Ich trinke einen kalten Kaffee.}{I drink a cold coffee.}
\end{itemize}
\end{exbox}
\end{grammarentry}

\section{Mini-Test}

\begin{grammarentry}{Übung}{exercise}
\GermanText{Setze die richtige Form von \textbf{gut} ein.}
\EnglishText{Insert the correct form of \textbf{gut}.}

\begin{enumerate}
    \item mit \underline{\hspace{2cm}} Deutsch
    \item mit dem \underline{\hspace{2cm}} Deutsch
    \item mit einem \underline{\hspace{2cm}} Beispiel
    \item ein \underline{\hspace{2cm}} Mann
    \item der \underline{\hspace{2cm}} Mann
    \item ohne \underline{\hspace{2cm}} Grund
\end{enumerate}

\GermanText{Lösungen: 1 gutem, 2 guten, 3 guten, 4 guter, 5 gute, 6 guten}
\EnglishText{Answers: 1 gutem, 2 guten, 3 guten, 4 guter, 5 gute, 6 guten}
\end{grammarentry}

\end{document}
